\chapter{双対空間}
    \section{基底変換と座標変換の関係}
        \begin{shadebox}
            $V$を$n$次元ベクトル空間とし($\realNumbers^n$や$\complexNumbers^n$とは限らないことに注意)、$\{\bm{v}_1,\dotsc,\bm{v}_n\}$と$\{{\bm{v}_1}',\dotsc,{\bm{v}_n}'\}$を$V$の基底とする。
            これらが$n$次正則行列$A$により次の関係式に従うとする。
            \[
                \begin{bmatrix}
                    {\bm{v}_1}' \\
                    \vdots \\
                    {\bm{v}_n}'
                \end{bmatrix}
                = A
                \begin{bmatrix}
                    \bm{v}_1 \\
                    \vdots \\
                    \bm{v}_n
                \end{bmatrix}
            \]
            $\{\bm{v}_1,\dotsc,\bm{v}_n\},\;\{{\bm{v}_1}',\dotsc,{\bm{v}_n}'\}$の双対基底をそれぞれ$\{\bm{v}^1,\dotsc,\bm{v}^n\}\,\{{\bm{v}^1}',\dotsc,{\bm{v}^n}'\}$とし、
            $\underline{\bm{x}} \coloneqq [x_1,\dotsc,x_n]^\top,\;\underline{\bm{x}}' \coloneqq [{x_1}',\dotsc,{x_n}']^\top$とする。
            このとき$[{\bm{v}^1}',\dotsc,{\bm{v}^n}']\underline{\bm{x}}' = [\bm{v}^1,\dotsc,\bm{v}^n]\underline{\bm{x}}$であれば$\underline{\bm{x}}' = A\underline{\bm{x}}$である。
            すなわち、双対空間での座標変換と元の空間の基底間の線形変換が同じ形になる。
        \end{shadebox}
        \begin{proof}
            \quad\par
            $[{\bm{v}^1}',\dotsc,{\bm{v}^n}']\underline{\bm{x}}' = [\bm{v}^1,\dotsc,\bm{v}^n]\underline{\bm{x}}$より
            \begin{align*}
                \underline{\bm{x}}' &= 
                \begin{bmatrix}
                    {\bm{v}_1}' \\
                    \vdots \\
                    {\bm{v}_n}'
                \end{bmatrix}
                [\bm{v}^1,\dotsc,\bm{v}^n]\underline{\bm{x}} = A
                \begin{bmatrix}
                    \bm{v}_1 \\
                    \vdots \\
                    \bm{v}_n
                \end{bmatrix}
                [\bm{v}^1,\dotsc,\bm{v}^n]\underline{\bm{x}} = A I_3\underline{\bm{x}} = A\underline{\bm{x}}
            \end{align*}
        \end{proof}